% ============================================================
% Chapter 5: Development
%   Section: News Acquisition: Discovery and Extraction
%     Subsection: The Extraction Cascade
% Included from chapters/05-development/02-news-acquisition/news-acquisition.tex
% ============================================================

\subsection{The Extraction Cascade}
\label{subsec:scraping-cascade}

\texttt{ExtractorService} (\path{services/scraper/extractor.py}) tries
three strategies in order of cost (Table~\ref{tab:cascade}).

\begin{table}[ht]
\centering
\small
\caption{The extraction cascade.}
\label{tab:cascade}
\begin{tabularx}{\linewidth}{@{}clX@{}}
\toprule
Step & Strategy & Cost \\
\midrule
1 & trafilatura           & One HTTP request. \\
2 & BeautifulSoup         & None: it parses the HTML step 1 already
                            fetched. \\
3 & Headless browser      & A second request, a full render with
                            JavaScript, seconds; never for evidence
                            pages. \\
\bottomrule
\end{tabularx}
\end{table}

\paragraph{When to escalate.} A failed step only passes the page on
when a different strategy could plausibly succeed
(Table~\ref{tab:escalation}). A dead source therefore costs one request
per URL, not one per strategy.

\begin{table}[ht]
\centering
\small
\caption{What happens after a failed step.}
\label{tab:escalation}
\begin{tabularx}{\linewidth}{@{}p{5.2cm}X@{}}
\toprule
Outcome of the step & Next \\
\midrule
No article in the page, too short, parse error
& The next strategy; an HTML parser gets the same page. \\
HTTP 401, 403 or 429 (a bot wall or a rate limit)
& Skip the HTML parsers, which would be refused the same way; go to the
  browser. \\
Any other HTTP error (404, 500\ldots), timeout, connection error,
refused by the guard
& Stop: no strategy can fix a page that is not there. \\
\bottomrule
\end{tabularx}
\end{table}

\paragraph{Step 1: trafilatura.} The main extractor, chosen because it
isolates the article body from navigation, comments and advertising on
most news layouts without any per-site configuration. On the
23~September~2026 check of two articles per readable source it
produced the article for every source that answered.

\paragraph{Step 2: BeautifulSoup.} Written for the layouts trafilatura
misses, it looks, in order, at:
\begin{enumerate}
    \item the CSS selectors in the source's own file, for layouts
    nothing else guesses;
    \item what the page declares about itself: the JSON-LD
    \texttt{NewsArticle}, whose \texttt{articleBody} often holds the
    full text behind a teaser, and the meta tags;
    \item the usual containers (\texttt{article},
    \texttt{[itemprop=articleBody]}, \texttt{main}), and finally
    whichever element holds the most paragraph text of its own.
\end{enumerate}
On the same pages it extracted comparable text, but on a CNN live blog
it picked a much smaller block, which is why it comes second.

\paragraph{Step 3: the headless browser.} Some pages build their text
with JavaScript, and some servers refuse clients that do not behave
like a browser. The third strategy renders the page with Playwright and
Chromium, and then reads the rendered HTML with the same two parsers,
so there is one notion of where an article is, not a third one written
against the browser's API. It is restricted in four ways:

\begin{itemize}
    \item \textbf{Never for evidence.} Evidence pages are read by the
    handful for every claim and already fall back to their search
    snippet, so rendering them would multiply the cost of the most
    expensive step by the most frequent purpose for very little. The
    browser is reserved for articles, ingestion, the enrichment page,
    the labelling batch, and the health checks that ask what ingestion
    would get.
    \item \textbf{Guarded like every other fetch.} The start URL is
    checked before a browser is launched, every request the page makes
    goes through the URL guard, and the redirect chain is checked after
    the page lands; content that arrived by way of a forbidden address
    is discarded.
    \item \textbf{Capped.} At most two renders run at once in the whole
    process (\path{BROWSER_MAX_CONCURRENCY}), one Chromium per
    render, because the browser's synchronous objects are bound to the
    thread that created them and extraction runs on many threads.
    \item \textbf{Optional.} The browser is an optional dependency,
    installed in the Docker image and usually absent on a development
    machine. Without it the attempt is \texttt{unavailable}, and the
    extraction keeps the previous step's outcome with a note that it
    was not escalated, so the statistics still say why the page failed.
\end{itemize}

It was checked live on 23~September~2026: a page whose text is built by
JavaScript has 29 characters of text in its HTML; trafilatura's result
is refused as too short, BeautifulSoup finds nothing, and the browser
renders 1,097 characters, for two requests and 11~seconds on a cold
start. An ordinary NASA page never gets past trafilatura: one request,
no browser.

\paragraph{Routing by source.} A URL posted for analysis arrives
without a source, so an El País article used to be treated as a
generic web page. The extractor now matches every URL against the
enabled sources by domain (the source's host or a subdomain of it, the
most specific match winning), so that the source's selectors and
browser flag apply and the stored article carries the real source.
Setting \texttt{requires\_javascript} puts the browser first, because
the two cheap steps are known to fail for that site and would each cost
a request to find that out again; they still run after it, for when no
browser is installed. The flag is meant to be set from evidence rather
than guessed: the monitoring page flags a domain whose every article
came from the browser. No source has needed it so far, and El País, the
case it was built for, answers 403 to the headless browser as well.

\paragraph{Validation.} An extraction is accepted only when its body
reaches the run's minimum length (500 characters by default, a per-run
threshold like every other tunable of the pipeline), which turns
paywall stubs, cookie walls and photo galleries into
\texttt{too\_short} rather than into articles with nothing to check.
